% ============================================================
% I. INTRODUCCIÓN
% ============================================================

\section{Introducción}

El muestreo, la cuantización y el análisis espectral son operaciones fundamentales al convertir una señal analógica en una representación digital procesable. El muestreo determina qué tan fielmente una señal continua puede reconstruirse a partir de sus muestras discretas; una tasa de muestreo insuficiente introduce \emph{aliasing}, un fenómeno en el que componentes de frecuencia alta se confunden con componentes de baja frecuencia. La cuantización, por su parte, aproxima cada muestra a uno de un número finito de niveles, introduciendo un error de cuantización inherente al proceso de conversión analógico-digital (DAC/ADC). Finalmente, la Transformada Rápida de Fourier (FFT) permite estudiar el contenido frecuencial de una señal ya digitalizada, útil tanto para caracterizar señales sintéticas como para diferenciar fuentes sonoras reales.

Este trabajo aborda estos conceptos en cuatro ejercicios independientes, implementados en el paquete \texttt{signal\_tools} del proyecto \texttt{WORK\_FOUR}: (1) el análisis de una señal muestreada real (\texttt{Muestra01.csv}); (2) el muestreo y cuantización de una señal seno de referencia a tres frecuencias de muestreo distintas, para ilustrar directamente el aliasing; (3) la comparación espectral de tres fuentes sonoras reales de naturaleza muy distinta (un instrumento de cuerda, un instrumento de percusión y un sonido animal); y (4) un análisis de la precisión temporal alcanzable al muestrear sensores analógicos en un microcontrolador ESP32-WROOM bajo tres mecanismos de concurrencia distintos, frente a una carga de cómputo que compite por el mismo procesador.

El objetivo es doble: por un lado, aplicar y verificar experimentalmente (mediante código y datos reales, no solo teoría) los conceptos de muestreo, aliasing y FFT; por otro, extender el análisis hacia un caso de ingeniería de sistemas embebidos donde la \emph{precisión temporal del muestreo} —no solo su frecuencia nominal— depende críticamente del mecanismo de concurrencia elegido.
